
% Issues in standard LLM genetation
Large Language Models (LLMs) are increasingly used for various tasks, including natural language generation~\cite{radford2019language} and code generation~\cite{chen2021evaluatinglargelanguagemodels}.
However, their outputs can suffer from issues such as hallucination~\cite{xu2024hallucinationinevitableinnatelimitation}, disclosure of private user information found in the training corpus \cite{wang2023decodingtrust}, as well as incorrect code generation in programming tasks.
When the output does not meet user expectations, users often have to restart the generation process with additional information in the prompt.
Alternatively, decoding strategies like beam search can generate multiple potential outputs for a single prompt, allowing for the selection of the most suitable response.
Both these approaches are computationally intensive and demand significant token generation, posing challenges in terms of efficiency and resource utilization.

% Issues in constrained LLM generation
%To address these issues, 
Recent techniques in context-free grammar (CFG) guided generation tried to address these issues by introducing constrained decoding techniques that ensure LLM outputs adhere to user-specified grammatical rules ~\cite{poesia2022synchromesh, willard2023efficient, guidance, geng2023grammar, ugare2024syncodellmgenerationgrammar, beurerkellner2024guiding}. 
These approaches typically involve various parsing techniques to analyze the LLM's partial outputs and determine the acceptable set of tokens {based on the defined grammar.}
% Issues in constrained LLM generation
While effective in producing \textit{syntactically} correct output, these techniques fall short of enforcing \textit{semantic} properties that extend beyond syntax.
%For example, ensuring that a variable name in LLM-generated code is defined before its use or that the generated text avoids harmful language, cannot be adequately captured by grammatical constraints alone. 
For example, grammatical constraints alone cannot adequately ensure that a variable name in LLM-generated code is defined before its use or that the generated text avoids harmful language. 

If an LLM generates a semantically incorrect output, the user typically must restart the generation from scratch.
Current grammar-guided generation tools fail to address this problem effectively, as they cannot detect semantic violations, or pause the generation at intermediate points. 
Additionally, navigation through the generation by naively backtracking a certain number of tokens from the end of the output to the part that caused the violation is very difficult. The main challenge is that the token-level abstraction provided by current LLM generation libraries~\cite{wolf-etal-2020-transformers, llamacpp} is not tied to the syntax of the underlying generation. 
%, making effective navigation through the generation difficult. %, making it often difficult to navigate through the generation effectively.
%
%Our key insight is that symbols in a grammar defined using Backus-Naur Form (BNF)  -- terminals (e.g., keywords, operators) and non-terminals (e.g., expressions, statements) -- offer a more intuitive and interpretable abstraction for navigating through the generation process. 
%
Our key insight is that symbols in a grammar, both terminals (e.g., keywords, operators) and non-terminals (e.g., expressions, statements)
offer a more intuitive and interpretable abstraction for navigating through \mbox{the generation process.}




% Our work
%To address these challenges, 
\paragraph{\Tool{}.} 
We introduce \Tool{}, a novel framework that provides a user-friendly interface for iteratively generating structured outputs from LLMs. 
Users specify a context-free grammar in the Backus-Naur Form (BNF) for the target output language, guiding the LLM to adhere to the syntax defined by the grammar.
%
Beyond syntax adherence, \Tool{} enables the user to programmatically check and correct for custom semantic properties of the generated output.
% For example, in a code generation task, the \Tool{} program can move forward and backward by a \textit{statement} or \textit{expression}, instead of a specific number of LLM tokens.
% This iterative control over the generation allows backtracking LLM generation and selectively resampling fragments of generation with any semantic violation. 
For example, in a code generation task, instead of moving forward or backward by a fixed number of LLM tokens, the \Tool{} program can navigate by higher-level abstractions such as \textit{statements} or \textit{expressions}.
This semantic-aware control enables selective resampling of fragments that violate desired properties, allowing for targeted corrections while preserving valid parts of the generation.
% \sasa{TBD: unclear last sentence; also For example sentence talks about syntax-directed generation, not that semantic analysis determines the move fwd/backward is needed. }


The key technical challenge to precise grammar-aware navigation is addressing
\textit{token misalignment} -- i.e., that LLM tokens from the model's fixed vocabulary do not directly correspond to lexical tokens associated with any specific grammar.
%
\Tool{} handles this issue by dynamically computing a mapping of grammar symbols to their corresponding positions in the partially parsed output. 
This capability enables efficient navigation both forward and backward through the generation process.
% Further, the generation efficiency is optimized by maintaining the state of the key-value (KV) cache. 
For each LLM generation task, \Tool{} maintains the history of generated tokens (as a tree of decoded tokens) that enables it to avoid regenerating the same tokens heuristically. 
\Tool{}'s intuitive interface can be used to program LLM generation algorithms that enhance specific semantic properties of the outputs by leveraging grammar symbols as navigational abstractions.

% Eval
Our evaluation presents three distinct scenarios, which demonstrate the effectiveness of \Tool{}. 
First, we illustrate how it can be used to improve the accuracy of LLM-generated SQL queries by enforcing additional semantic constraints.
\Tool{} achieves 18.5\% mean improvement over the state-of-the-art grammar-guided generation technique~\cite{ugare2024syncodellmgenerationgrammar}.
Second, we show how \Tool{} effectively reduces privacy leaks in LLM-generated text from 51.4\% to 0\%, thus successfully safeguarding sensitive information while maintaining the quality of response.
Third, we show that \Tool{} improves the accuracy of LLM-generated Vega-lite specification (a subset of JSON for data visualization) by 17.8\% by enforcing semantic constraints. 

\noindent{\bf Contributions.} The main contributions of this paper are:

\vspace{-.06in} \begin{itemize}[leftmargin=*]
\itemsep 1pt 
\parskip 2pt 
%
\item We present \Tool{}, the first framework that uses grammar symbols as abstractions for navigating LLM generation both forward and backward. 

\item We introduce an algorithm that enables efficient and accurate control of the LLM generation through grammar symbol abstraction by maintaining the decoding history and the \mbox{LLM key-value cache.}

\item We demonstrate how \Tool{} enhances specific semantic properties in LLM-generated outputs through three scenarios, addressing issues of privacy leaks and accuracy in SQL \add{and Vega-Lite specification generation}. \end{itemize}

